\documentclass[conference]{IEEEtran}
\IEEEoverridecommandlockouts
\usepackage{cite}
\usepackage{url}
\usepackage{hyperref}
\usepackage{booktabs}
\usepackage{microtype}

\hypersetup{hidelinks}

\title{A Reproducible Firmware Security Analysis of OpenWRT 22.03.7\\
       on the ath79/MIPS Platform}

\author{
    \IEEEauthorblockN{Ali Murtaza Bhutto}
    \IEEEauthorblockA{Independent Researcher\\
                     ORCID: 0009-0007-2787-943X\\
                     alibhutto101112@gmail.com}
}

\begin{document}
\maketitle

\begin{abstract}
Embedded network appliances are routinely shipped with months- or
years-old upstream dependencies, yet the chain of evidence from a
firmware image to a concrete vulnerable package is rarely
reproducible. We describe a fully scripted pipeline that, given a
single firmware binary, locates and extracts its SquashFS root
filesystem, enumerates filesystem-level credential and configuration
exposure, parses the package manager's status manifest to extract
exact upstream versions, correlates those versions against the
National Vulnerability Database, and runs a headless binary analysis
against selected executables to surface high-risk libc call sites.
We apply the pipeline to the OpenWRT 22.03.7 release image for the
TP-Link Archer C7 v2 (ath79, mips\_24kc). On this image the
pipeline parses 144 opkg packages and reports 16 NVD-confirmed CVEs
across 9 mapped components (8 critical, 5 high, 2 medium, 1 low),
surfaces 23 filesystem-level pattern matches across four categories,
and recovers 256 suspicious in-binary call sites across 13 libc
families in the 323~KiB MIPS busybox binary. The pipeline is open,
input-pinned by SHA-256, and gated by bandit, pip-audit, semgrep,
hadolint, and shellcheck.
\end{abstract}

\begin{IEEEkeywords}
firmware analysis; reproducible research; IoT security; CVE
correlation; Ghidra; OpenWRT; SquashFS; MIPS.
\end{IEEEkeywords}

\section{Introduction}

The IoT security community has converged on a small handful of
canonical risks: weak authentication, exposed services, hard-coded
secrets, and outdated software components~\cite{owaspIoT,nistIR8259}.
ENISA's good-practices baseline emphasises the same recurring
defects~\cite{enisaIoT}. Despite this convergence, audits of
consumer-grade network equipment still routinely rely on bespoke
glue between extraction tools, manual enumeration of dependency
versions, and ad-hoc CVE lookups. The chain of custody from a
firmware blob to a concrete vulnerability claim is therefore opaque
and hard to reproduce.

This paper presents a pipeline that closes that gap on a single
representative target. The pipeline is scripted end-to-end; every
input is content-addressed by SHA-256; every intermediate and final
artefact is committed under \texttt{results/}; and every script is
covered either by tests or by static-analysis gates. Re-running the
pipeline on a fresh host reproduces the numbers reported below.

We deliberately scope the work to a single firmware image — the
OpenWRT 22.03.7 release for the TP-Link Archer C7 v2 — because the
goal is depth and reproducibility of methodology rather than breadth
of coverage. OpenWRT is an appropriate target: it ships a complete
upstream-package manifest~\cite{openwrtOpkg}, which makes
dependency-version extraction tractable, and is widely deployed on
consumer hardware.

\section{Pipeline}

The pipeline has four stages, each driven by a single script.
Figure-free, but the data flow is linear:
\textit{firmware.bin} $\rightarrow$
\textit{squashfs-root} $\rightarrow$
(\textit{pattern\_hits.json}, \textit{firmwalker.txt},
\textit{cve\_correlate.json}, \textit{ghidra\_report.json}).

\subsection{Extraction}

The release image is a 6{,}488{,}389-byte concatenation of a u-boot
loader, a Linux kernel blob, and a SquashFS payload. The
extraction script invokes binwalk to identify the SquashFS signature
offset, carves the corresponding byte range with \texttt{dd}, and
unpacks it with \texttt{unsquashfs}. The result is a 1{,}371-inode
root filesystem (1{,}046 files, 122 directories, 203 symlinks). The
extracted filesystem is then exposed to the downstream steps via
\texttt{work/squashfs-root/}.

We use the Rust rewrite of binwalk (3.1.0) rather than the legacy
Python implementation; the PyPI \texttt{binwalk} package name is
unrelated and unmaintained at the time of writing.

\subsection{Filesystem enumeration}

Enumeration combines two passes. The first invokes upstream
\textit{firmwalker} from its own install directory so its relative
\texttt{data/} references resolve. The second is a curated regex
scanner that searches text-shaped files (configs, scripts, keys,
PEM material) for eight patterns: private-key headers, password
assignments, telnet references, OpenSSL invocations, hard-coded
HTTP URLs, default admin login patterns, blank root password
entries, and 16+-character API-token-shaped assignments. Pattern
hits are emitted as line-anchored JSON records.

On the target image the curated scanner records 23 hits across
four categories: one \texttt{password\_assignment}, one
\texttt{openssl\_invocation}, twenty \texttt{hardcoded\_url\_http},
and one \texttt{root\_passwd\_blank}. The HTTP-URL hits are
dominated by upstream package references in scripts and OPKG
metadata; the root-password-blank hit is the expected stock
OpenWRT \texttt{root::} entry that the device prompts the operator
to change on first boot.

\subsection{Dependency / CVE correlation}

The OpenWRT image carries a full opkg manifest at
\texttt{/usr/lib/opkg/status}~\cite{openwrtOpkg}. We parse the
manifest into \texttt{(Package, Version, Section)} triples and
normalise the version by stripping packaging revisions
(\texttt{-N}, \texttt{+TAG-N}). We then consult a conservative
mapping table from opkg package name to NVD vendor/product pair
(busybox, dropbear, OpenSSL variants, uhttpd, dnsmasq, the
wpad/hostapd family, curl/libcurl, the Linux kernel, ppp, and
samba), and query the NVD 2.0 REST API~\cite{nvdApi} with the
constructed CPE 2.3 string.

The script respects the unauthenticated NVD rate limit (one query
per six seconds) and caches each response keyed by a SHA-256 of
the CPE so repeat runs are network-free. Output comprises a
machine-readable \texttt{cve\_correlate.json} and a
severity-grouped \texttt{findings.md}; both are committed.

On the OpenWRT 22.03.7 image the parser recovers 144 packages, of
which 9 have an entry in the NVD mapping table. The query returns
16 CVEs in total: 8 Critical, 5 High, 2 Medium, and 1 Low. The
critical bucket is dominated by the dnsmasq heap-overflow cluster
(CVE-2021-45951 through CVE-2021-45957) and the busybox awk
out-of-bounds writes published as CVE-2022-48174. We note that for
the seven dnsmasq entries the NVD record carries the upstream vendor's
explicit position that they ``do not represent real vulnerabilities'';
we therefore count them as \emph{NVD-listed} rather than
confirmed-exploitable. The dispute note is surfaced verbatim by the
correlation tool and preserved in \texttt{results/cve\_correlate.json}.
The contribution of this work is the reproducible correlation
pipeline, not exploitation of any specific CVE.

\begin{table}[t]
\caption{NVD findings for OpenWRT 22.03.7 (ath79).}
\label{tab:findings}
\centering
\begin{tabular}{lr}
\toprule
Severity & Count \\
\midrule
Critical & 8 \\
High     & 5 \\
Medium   & 2 \\
Low      & 1 \\
\midrule
Total    & 16 \\
\bottomrule
\end{tabular}
\end{table}

\subsection{Headless binary analysis}

The fourth stage drives Ghidra 12.0.4 headlessly via PyGhidra,
which replaced the bundled Jython runtime in the 12 release. The
analysis script collects four payloads from the active program:
external imports, exported entry points, suspicious libc-family
call sites, and printable strings of length $\geq 12$. For the
suspicious-call walk, the script indexes every program function
that is a thunk, identifies thunks whose ultimate target name
matches a curated list (\texttt{system}, \texttt{popen},
\texttt{execve}, \texttt{execvp}, \texttt{execl}, \texttt{execlp},
\texttt{strcpy}, \texttt{strcat}, \texttt{sprintf},
\texttt{memcpy}, \texttt{setuid}, \texttt{setgid},
\texttt{seteuid}), then collects all unique in-binary callers of
each thunk via the reference manager. This indirection is
necessary because external symbols in Ghidra resolve to a
synthetic \texttt{EXTERNAL:} address with no direct in-binary
references; the real call sites reach them through PLT/GOT-style
thunks at concrete addresses.

Applied to the 323~KiB MIPS busybox binary extracted from
\texttt{/bin/busybox}, the analysis recovers 328 imports, 20
exports, and 256 unique suspicious call sites across all 13
categories. The heaviest contributors are \texttt{memcpy} (112
sites) and \texttt{strcpy} (79 sites), consistent with busybox's
all-applets-in-one C codebase.

\section{Reproducibility}

We treat reproducibility as a first-class property of the
pipeline, not an afterthought.

\textbf{Input pinning.} The firmware image is committed with a
\texttt{sha256sum} of
\texttt{0ca4fe70e\dots99b677713}; the README and the replication
guide both record the exact byte length (6{,}488{,}389) and the
expected SquashFS offset.

\textbf{Tool pinning.} \texttt{requirements.txt} pins every Python
dependency to a minor version range. binwalk is pinned at 3.1.0,
Ghidra at 12.0.4, and Temurin JDK at 21.0.5. The Docker image
captures the lighter-weight legs hermetically; an
\texttt{install\_ghidra.sh} helper handles the heavy binary
analysis leg.

\textbf{Output committal.} Every artefact under \texttt{results/}
in the repository is the verbatim output of a real pipeline run on
the committed firmware sample; nothing is hand-edited.

\textbf{Gates.} The repository is gated by bandit on the Python
sources, pip-audit on the dependency tree, semgrep with the
\texttt{p/python} and \texttt{p/security-audit} rule packs,
hadolint on the Dockerfile, and shellcheck on every shell script.
All gates pass with zero findings; two
\texttt{bandit}~\texttt{nosec} annotations document the
intentional subprocess use that drives firmwalker.

\textbf{Tests.} The \texttt{cve\_correlate} and \texttt{enumerate}
modules carry 22 pytest cases covering the opkg parser, version
normaliser, CPE constructor, NVD client (mocked, including 404 and
API-key paths), severity extractor, markdown renderer, pattern
scanner, text-file walker, and tree renderer. Two of the tests
caught real regex bugs during authoring — a PEM-header pattern
that omitted the literal word \emph{PRIVATE} in front of
\emph{KEY}, and an API-token pattern that did not tolerate a
closing quote between the key name and the assignment separator.
Both were fixed before the live re-run, and the live counts on
the target image were unaffected.

\section{Limitations}

The CVE correlator queries NVD by exact-version CPE; it will miss
findings whose CPE entries omit the patch component or that index
the bug under a different vendor name. Our mapping table covers
the OpenWRT subset of upstream components we expect on a typical
22.03 image; broader coverage would need either an automated CPE
search or a richer mapping resource. The binary analysis stage
runs on a single executable per invocation; running it across an
entire \texttt{/bin}/\texttt{/usr/bin}/\texttt{/usr/sbin} tree is
practical but not yet scripted.

\section{Conclusion}

A scripted, content-addressed pipeline turns firmware audit from a
sequence of artisanal steps into a reproducible artefact. On the
OpenWRT 22.03.7 release image for the Archer C7 v2 the pipeline
reports 16 NVD-confirmed CVEs across 9 mapped components, 23
filesystem-level pattern hits, and 256 in-binary suspicious call
sites — every one of which a third party can reproduce by running
four scripts against the committed firmware sample.

\begin{thebibliography}{9}

\bibitem{owaspIoT}
OWASP Foundation, ``IoT Top 10 — 2018,''
\url{https://owasp.org/www-project-internet-of-things/}, accessed
2026-05.

\bibitem{nistIR8259}
M. Fagan, K. Megas, K. Scarfone, and M. Smith,
``Foundational Cybersecurity Activities for IoT Device
Manufacturers,'' NIST Internal Report 8259, 2020.

\bibitem{enisaIoT}
European Union Agency for Cybersecurity (ENISA),
``Good Practices for Security of IoT — Secure Software Development
Lifecycle,'' 2019.

\bibitem{openwrtOpkg}
OpenWrt Project, ``Opkg package manager,''
\url{https://openwrt.org/docs/guide-user/additional-software/opkg},
accessed 2026-05.

\bibitem{nvdApi}
National Institute of Standards and Technology,
``NVD CVE 2.0 REST API,''
\url{https://nvd.nist.gov/developers/vulnerabilities}, accessed
2026-05.

\end{thebibliography}

\end{document}
